\newpage
\section{\tkzClass{parallelogram} 类}
变量 \tkzVar{parallelogram}{P} 持有一个表，用于存放平行四边形。它是可选的，你可以自由选择变量名。不过，为清晰和一致起见，推荐使用 \code{P}。如果你使用自定义变量（例如 parall），则必须手动初始化它。如果使用了 \code{P}，\code{init\_elements()} 函数会重新初始化该表。


\subsection{创建平行四边形}
\label{sub:creating_a_parallelogram}

\tkzClass{parallelogram} 类根据它的四个顶点创建一个平行四边形。要由三个点计算第四个顶点，使用 \tkzFct{parallelogram}{parallelogram.fourth(za, zb, zc)}。

\medskip
所得对象存放在 \tkzVar{parallelogram}{P} 中。你可以自由使用别的名字，但为一致起见推荐使用 \tkzVar{parallelogram}{P}。

\begin{mybox}
\begin{verbatim}
P.ABCD = parallelogram:new(z.A, z.B, z.C, z.D)
\end{verbatim}
\end{mybox}

\textbf{简写形式。}

你也可以使用简写形式：

\begin{mybox}
\begin{verbatim}
P.ABCD = parallelogram(z.A, z.B, z.C, z.D)
\end{verbatim}
\end{mybox}

\subsection{平行四边形的属性}

各点按正向（逆时针）给出。有一个检验会检查两条对角线是否具有相同的中点；否则会引发错误。

\begin{mybox}
创建 | P.new = parallelogram(z.A,z.B,z.C,z.D)|
\end{mybox}


  \bgroup
  \catcode`_=12
  \small
  \captionof{table}{平行四边形的属性。}\label{parallelogram:attributes}
  \begin{tabular}{lll}
  \toprule
  \textbf{属性}         & \textbf{用法}  &  \\
  \midrule
  \tkzAttr{parallelogram}{pa}   & |z.A = P.new.pa|     &  \\
  \tkzAttr{parallelogram}{pb}   & |z.B = P.new.pb|     &  \\
  \tkzAttr{parallelogram}{pc}   & |z.C = P.new.pc|     &  \\
  \tkzAttr{parallelogram}{pd}   & |z.D = P.new.pd|     &  \\
  \tkzAttr{parallelogram}{type} & |P.new.type= 'parallelogram'|&  \\
  \tkzAttr{parallelogram}{center}    & |z.I = P.new.center|  & 对角线交点 \\
  \tkzAttr{parallelogram}{ab}   & |P.new.ab|    &  过两个顶点的直线 \\
  \tkzAttr{parallelogram}{ac}   & |P.new.ac|           &  同上。 \\
  \tkzAttr{parallelogram}{da}   & |P.new.da|           &  同上。 \\
  \tkzAttr{parallelogram}{bc}   & |P.new.bc|           &  同上。 \\
  \tkzAttr{parallelogram}{bd}   & |P.new.bd|           &  同上。 \\
  \tkzAttr{parallelogram}{cd}   & |P.new.cd|           &  同上。 \\
  \bottomrule %
  \end{tabular}
  \egroup


\subsubsection{示例：属性}
\label{ssub:example_attributes}
\begin{minipage}{.5\textwidth}
\begin{verbatim}
\directlua{
 init_elements()
 z.A = point(0, 0)
 z.B = point(4, 1)
 z.C = point(7, 5)
 z.D = point(3, 4)
 P.new = parallelogram(z.A, z.B, z.C, z.D)
 z.B = P.new.pb
 z.C = P.new.pc
 z.D = P.new.pd
 z.I = P.new.center}
\begin{tikzpicture}
\tkzGetNodes
\tkzDrawPolygon(A,B,C,D)
\tkzDrawPoints(A,B,C,D)
\tkzLabelPoints(A,B)
\tkzLabelPoints[above](C,D)
\tkzDrawPoints[red](I)
\end{tikzpicture}
\end{verbatim}
\end{minipage}
\begin{minipage}{.5\textwidth}
\directlua{
 init_elements()
 z.A = point(0, 0)
 z.B = point(4, 1)
 z.C = point(7, 5)
 z.D = point(3, 4)
 P.new = parallelogram(z.A, z.B, z.C, z.D)
 z.B = P.new.pb
 z.C = P.new.pc
 z.D = P.new.pd
 z.I = P.new.center}

\begin{center}
  \begin{tikzpicture}
  \tkzGetNodes
  \tkzDrawPolygon(A,B,C,D)
  \tkzDrawPoints(A,B,C,D)
  \tkzLabelPoints(A,B)
  \tkzLabelPoints[above](C,D)
  \tkzDrawPoints[red](I)
  \end{tikzpicture}
\end{center}
\end{minipage}

\newpage
\subsection{平行四边形的函数}

\begin{center}
  \bgroup
  \catcode`_=12
  \small
  \captionof{table}{平行四边形的函数。}\label{parallelogram:methods}
  \begin{tabular}{ll}
  \toprule
  \textbf{函数} & \texttt{参考}    \\
  \midrule
  \tkzMeth{parallelogram}{new (za, zb, zc, zd)} & \ref{ssub:example_attributes}\\
  \tkzFct{parallelogram}{parallelogram.fourth (za,zb,zc)} & \ref{ssub:parallelogram_with_fourth_method} \\
  \tkzMeth{parallelogram}{get(<i>)} & [\ref{ssub:method_parallelogram_get}]\\
  \bottomrule %
  \end{tabular}
  \egroup
\end{center}

\subsubsection{方法 \tkzMeth{parallelogram}{get(<i>)}}
\label{ssub:method_parallelogram_get}

该方法返回定义平行四边形对象的各个点。

\begin{itemize}
  \item \code{P.ABCD:get()} 按顺序返回四个顶点：\code{pa}、\code{pb}、\code{pc}、\code{pd}；
  \item \code{P.ABCD:get(i)}（\code{i} 取 1 到 4）只返回第 \code{i} 个顶点；例如 \code{P.ABCD:get(2)} 等价于 \code{P.ABCD.pb}。
\end{itemize}

\begin{tkzexample}[code only]
\directlua{
  init_elements()
  z.A = point(0, 0)
  z.B = point(2, 0)
  z.C = point(3, 2)
  P.ABCD = parallelogram.fourth(z.A, z.B, z.C)
  z.E, z.F, z.G, z.H = P.ABCD:get()
  z.K = P.ABCD:get(3)
}
\end{tkzexample}

\subsubsection{方法 \tkzMeth{parallelogram}{new(pt,pt,pt,pt)}}
\label{ssub:function_parallelogram_new}

\begin{tkzexample}[latex=.5\textwidth]
\directlua{
 z.A = point(0, 0)
 z.B = point(4, 1)
 z.C = point(7, 5)
 z.D = point(3, 4)
 P.ABCD = parallelogram(z.A, z.B, z.C, z.D)
 z.B = P.ABCD.pb
 z.C = P.ABCD.pc
 z.D = P.ABCD.pd
 z.I = P.ABCD.center}
\begin{tikzpicture}
\tkzGetNodes
\tkzDrawPolygon(A,B,C,D)
\tkzDrawPoints(A,B,C,D)
\tkzLabelPoints(A,B)
\tkzLabelPoints[above](C,D)
\tkzDrawPoints[red](I)
\end{tikzpicture}
\end{tkzexample}

\subsubsection{方法 \tkzMeth{parallelogram}{fourth(pt,pt,pt)}}
\label{ssub:parallelogram_with_fourth_method}
用平行四边形补全一个三角形（见下一个示例）

\begin{minipage}{.5\textwidth}
\begin{verbatim}
\directlua{
init_elements()
z.A = point(0, 0)
z.B = point(3, 1)
z.C = point(4, 3)
P.four= parallelogram.fourth(z.A, z.B, z.C)
z.D = P.four.pd
z.I = P.four.center}
\begin{tikzpicture}[ scale = .75]
\tkzGetNodes
\tkzDrawPolygon(A,B,C,D)
\tkzDrawPoints(A,B,C,D)
\tkzLabelPoints(A,B)
\tkzLabelPoints[above](C,D)
\tkzDrawPoints[red](I)
\end{tikzpicture}
\end{verbatim}
\end{minipage}
\begin{minipage}{.5\textwidth}
\directlua{
 init_elements()
 z.A = point(0, 0)
 z.B = point(3, 1)
 z.C = point(4, 3)
 P.four= parallelogram.fourth(z.A, z.B, z.C)
 z.D = P.four.pd
 z.I = P.four.center}
\begin{center}
  \begin{tikzpicture}[scale = .75]
  \tkzGetNodes
  \tkzDrawPolygon(A,B,C,D)
  \tkzDrawPoints(A,B,C,D)
  \tkzLabelPoints(A,B)
  \tkzLabelPoints[above](C,D)
  \tkzDrawPoints[red](I)
  \end{tikzpicture}
\end{center}
\end{minipage}
\endinput